\scanpage{110}
Récapitulons le formulaire pour
\diagram{\begin{tikzcd}[column sep=large]
 k\arrow[r,"\deleted{\unsure{\alpha}}\,\alpha"] & A\ (k.\mathrm{alg})\arrow[l,bend left=35,"t"]\arrow[l,bend left=65,"n"]
\end{tikzcd}}
\margin{$\alpha$ hom d'anneaux comm. $\alpha(\lambda)=\lambda.1_A$.
On pose $\tau(x)=\alpha t(x)$, $\unsure{\nu|\gamma}(x)=\alpha n(x)$.}

\textcircled{1} $t$ \unsure{est} $k$-linéaire, $\boxed{t(1_A)=2}$.

\textcircled{2} $x\longmapsto\bar x\overset{\mathrm{def}}{=}\tau(x)-x:A\longrightarrow A$
\margin{(autom. involutif du $k$-\uline{module} $A$, avec $\bar1_A=1_A$)}
est un \uline{automorphisme} d'alg., i.e. on a
\[
 \boxed{\overline{xy}=\bar x\bar y}\qquad\text{i.e.}\quad
 2xy=(\tau(xy)-\tau(x)\tau(y))+\tau(y)x+\tau(x)y
\]

\textcircled{3}
\[
 \boxed{\left\{\begin{aligned}n(\lambda x)&=\lambda^2n(x)\\n(x+y)&=n(x)+n(y)+t(x\bar y)\end{aligned}\right.}
 \qquad\lambda\in k,\ x\in A
\]
i.e. $n$ est quadratique, la forme bil. ass. est \unsure{ici}
\[
 \varphi_n(x,y)=t(x\bar y)
\]

\textcircled{4} $\boxed{n(1)=1}$
\margin{(\uline{NB} compte \deleted{\unsure{dans}} tenu de la formule d'additivité dans $3^\circ$, cela implique $t(1)=2$)}

\textcircled{5} $\boxed{\alpha(n(x))=x\bar x}$\qquad$(x\in A)$
\margin{qui implique $3^\circ,4^\circ$ si $\alpha$ injectif}

\textcircled{6} (\uline{facultatif}), \deleted{\unsure{si}}\inserted{et} (si $k$ loc. libre de rang $n$ sur $k_0$, et $A$ loc. libre de rang $2n$ sur $k_0$)
\[
 \boxed{\begin{aligned}
 \operatorname{Tr}_{k/k_0}(t(x))&=\operatorname{Tr}_{A/k_0}(x)\\
 N_{k/k_0}(n(x))&=N_{A/k_0}(x)
 \end{aligned}}\qquad x\in A
\]
La donnée de \inserted{$A$ sur $k$ avec} $\alpha,t,n$ satisfaisant $1^\circ$ à $5^\circ$ sera appelée algèbre \uline{quasi-quadratique} sur $k$.
On va maintenant \unsure{poursuivre l'étude de la} situation \deleted{\unsure{bi}}\inserted{quasi} \uline{quadratique}
\mspage{109}
